\subsection{data parallel 剩下的问题}

即使 ZeRO 消除了大部分重复状态，data parallel 仍受两个边界约束。第一，global batch 不能无限增大，否则优化质量和样本效率会变化；第二，ZeRO-1/2 并不切分模型参数本身，也不必然解决 activation memory，因此超大模型仍需要沿模型内部维度继续切分。

\begin{figure}[H]
\centering
\includegraphics[width=0.86\textwidth]{slide-029.jpg}
\caption{Slide 29：data parallel 的 compute scaling 受 batch size 限制，机器数不能无脑超过有效 batch。}
\figsource{CS336 Spring 2026 Lecture 8 官方幻灯片。}
\end{figure}

\begin{knowledgebox}{读图：Slide 29 说的是“吞吐扩展到哪里停”}
data parallel 增加 rank 会增加 global batch 或减小 per-rank batch。若 global batch 过大，优化收益递减；若 per-rank batch 过小，GPU 利用率下降且通信占比上升。因此 DP 不是无限扩展 compute 的答案。
\end{knowledgebox}

\begin{figure}[H]
\centering
\includegraphics[width=0.86\textwidth]{slide-030.jpg}
\caption{Slide 30：ZeRO-1/2 不解决模型本身放不下和 activation memory，仍需 model parallelism。}
\figsource{CS336 Spring 2026 Lecture 8 官方幻灯片。}
\end{figure}

\begin{warningbox}{读图：Slide 30 把 ZeRO 的边界讲清楚}
ZeRO-1/2 主要分片 optimizer state 和 gradients，参数仍复制。ZeRO-3 能分片参数，但 activation memory 仍然随 sequence length、batch、层数增长。若模型层或激活本身太大，就要切模型内部结构，而不是只靠 data parallel 的状态分片。
\end{warningbox}


\singleslide{slides-images/slide-031.jpg}{从 data-state sharding 转向 model parallel：切参数后通信 activations，而不是临时传参数。}{31}

Slide 31 给出 ZeRO-3 与 model parallel 的分界。两者都把参数分散到设备，但 ZeRO-3 在模块执行前重新拼完整参数；pipeline/tensor/expert parallel 让每张卡永久计算模型的一部分，并在边界传 activation。选择依据不是“都能省参数显存”，而是参数通信与 activation 通信哪个更适合当前网络和 batch。

\begin{knowledgebox}{讲义补充：源 Slide 31 的 model parallel 和 ZeRO-3 有何不同}
两者都可能把参数放在多张卡上，但通信语义不同。ZeRO-3 保持 data-parallel 视角，需要时 gather 参数；model parallel 则让不同 rank 持有模型不同结构部分，并在前向/反向中传递激活或 partial results。它改变了模型计算图的并行执行方式。
\end{knowledgebox}

